\section{Introduction}
You are looking for a parking spot, driving toward your destination. Each spot
you pass is either free or taken, and once you drive past a free one you cannot
return. Stop too early and you walk a long way; hold out for the perfect spot
and you may overshoot and find nothing. When should you take what you have?

The popular answer invokes the \emph{secretary problem}: examine and reject the
first $37\%$ of the options, then take the next one that beats everything seen
so far. The cutoff $1/e\approx\OneOverE$ is genuinely optimal --- but only for
the exact problem the secretary story describes: a fixed, known number of
candidates, no recall, no cost to looking, and a payoff that rewards \emph{only}
selecting the single best. The parking problem violates several of these
assumptions, and the $37\%$ answer is simply wrong for it.

\paragraph{Question.} Two questions organise this paper. (i) How fragile is the
$37\%$ rule --- how far does the optimal stopping threshold move as we relax the
secretary assumptions toward the parking problem and toward costly search?
(ii) If an agent is given only the reward signal and no theory, does it
\emph{rediscover} the correct threshold from experience?

\paragraph{Approach.} We cast the classic secretary problem, the parking
problem, and two variants (a per-observation cost; cardinal ``full
information'') inside one MDP skeleton: observe, then \texttt{stop} or
\texttt{continue}. For the secretary problem the optimal policy has a
closed form; for the parking problem we obtain exact ground truth by backward
induction. Against these references we run tabular $Q$-learning and a
one-parameter policy-gradient (REINFORCE) learner.

\paragraph{Findings.} (i) The $37\%$ rule emerges cleanly only in the classic
rank-only, cost-free limit: at $n=\SecNhi$ the optimum sits at
$x^*=\SecXstarHi$ with success $\SecPstarHi$, but for $n=\SecNlo$ it is already
$\SecXstarLo$. (ii) Transplanted onto parking, the $37\%$ rule is disastrous:
regret up to $\ParkNaiveRegretHi\%$ at occupancy $p=\ParkPhi$. (iii) A search
cost of $\CostMax$ per look pulls the optimal secretary cutoff down to
$\CostMaxXstar$. (iv) $Q$-learning recovers the secretary optimum
($\AgentQregret\%$ regret at $n=\AgentSecN$) and the parking optimum (exact at
$p=\ParkPhi$) from reward alone; the one-parameter policy-gradient agent finds
the right regime but with higher variance.

\paragraph{Contributions.} A unified MDP treatment of the stopping family with
exact-DP ground truth; a quantification of how badly the $37\%$ rule fails off
its home turf; and an empirical study of what model-free agents recover.
